\documentclass[18pt]{article}
\usepackage{geometry}
\geometry{top=2cm}
\usepackage{amsmath,amsfonts,amssymb}
\usepackage{natbib}
\usepackage{booktabs}
\usepackage[utf8]{inputenc}
\usepackage[T1]{fontenc}
\usepackage{amssymb}
\usepackage[version=4]{mhchem}
\usepackage{stmaryrd}
\usepackage{bbold}
\usepackage{nopageno}
\usepackage{hyperref}
\usepackage{bookmark}
\author{Adam Handwerger}
\pagestyle{plain}

\begin{document}
\begin{flushleft}
1 A Nice Kernel. Let $\mathcal{X}=\mathbb{R}^d$. Define:\\
\vspace{.4cm}
\hspace{3cm} $k(x_k,x_l)=\sum\limits_{i,\lambda_i>0} \lambda_i \phi_i(x_k)\phi(x_l)$\\
\vspace{.4cm}
where $\lambda_i$ are the eigenvalues and $\phi_i(x)$ are the eigenfunctions of K.\\
\vspace{.4cm}
1. Verify that K is pos. definite.\\
\vspace{.4cm}
Let $\alpha \in \mathbb{R}^n$, then\\
\vspace{.4cm}
$\sum\limits_{kl} \alpha_k \alpha_l  K(x_k,x_l)=\sum\limits_{kl} \sum\limits_i \alpha_k \alpha_l \lambda_i \phi_i(x_k) \phi_i(x_l)=$\\
\vspace{.2cm}
$\sum\limits_i \left(\sum\limits_k \sqrt{\lambda_i}\alpha_k \phi(x_k) \sum\limits_l \sqrt{\lambda_i} \alpha_l \phi(x_i)\right)=$\\
\vspace{.4cm}
$\sum\limits_i (\varPhi_i(x))^2 \geq 0,\text{ where} \hspace{.2cm}\varPhi_i(x) \equiv \sum\limits_m \sqrt{\lambda_i}\alpha_m \phi(x_m)$\\
\vspace{.4cm}
*fixed missing $\alpha_i$s from first version.\\
\vspace{.4cm}
2. Show that the Hilbert space H is RKHS.\\
\vspace{.4cm}
K is in H since\\
\vspace{.4cm}
$f_x(\cdot)=K(x,\cdot)=\sum\limits_k \gamma_k K(x_k,\cdot)=\sum\limits_k \gamma_k \sum\limits_i \lambda_i \phi_i(x_k) \phi_i(\cdot)=$\\
$\sum\limits_i \lambda_i \left(\sum\limits_k \gamma_k \phi_i(x_k)\right) \phi_i(\cdot)$\\
\vspace{.4cm}
Define $\beta_i=\lambda_i \left(\sum\limits_k \gamma_k \phi_i(x_k)\right)$, then\\
\vspace{.4cm}
$f_x(\cdot)=K(x,\cdot)=\sum\limits_i \beta_i \phi_i(\cdot) \implies K(x_i,x_j)=\sum\limits_i \beta_i \phi_i(x_j)=f_x(x_j)$\\
\vspace{.4cm}
Thus K has the form of an $f \in H$.\\
\vspace{.4cm}
3. K has the reproducing property since\\
\vspace{.4cm}
$\langle f_\alpha, K(x,\cdot)\rangle=\langle \sum\limits_i \alpha_i \phi_i(\cdot),\sum\limits_i \lambda_i \phi_i(x)\phi_i(\cdot)\rangle=$\\
$\sum\limits_i \alpha_i \phi_i(x) \langle \phi_i(\cdot),\phi_i(\cdot) \rangle=\sum\limits_i \alpha_i \phi(x)=f_\alpha(x)$\\
\vspace{.4cm}
This follows since eigenfuctions are orthonormal. Thus H is RKHS.\\
\vspace{.4cm}
4. Show that the inner product in H can be written as $\sum_i \alpha_i \beta_i$\\
\vspace{.4cm}
Take $\alpha^{\prime}_i=\alpha_i/\sqrt{\lambda_i}$. And redefine H with $\alpha\rightarrow \alpha^{\prime}$ and inner product:\\
\vspace{.4cm}
$\sum\limits_i \dfrac{\alpha_i \beta_i}{\lambda_i}=\sum\limits_i\dfrac{ \sqrt{\lambda_i} \alpha_i^{\prime} \sqrt{\lambda_i} \beta_i^{\prime}}{\lambda_i}=\sum_i \alpha_i^{\prime} \beta_i^{\prime}$\\
\vspace{.4cm}
*Added forgotten primes in first versions





















\end{flushleft}
\end{document}
